%=============================================================================
% Wave 4, Iteration 10: Materials Science Update
% AGENT: MATERIALS SCIENCE (W4i10) | Model: Opus 4.6
%
% INHERITED FROM W3i7 MatSci:
%   - Passive FOM inverts active: Nb dominates passive, YBCO worst
%   - Si3N4 MEMS: Q>=10^9 demonstrated (room T), 1.2 ng, 620 kHz
%   - Complete fab spec (Design C): 500 um defect, 50 nm thick
%   - Nb recipe: RRR>=300, N-doped, Q0>2e11 at 20mK
%   - Superfluid He-3: Track 3 (gravitational probe, not lambda constraint)
%   - Novel materials (nanowires, topological, BEC): all killed
%   - ng MEMS at 1/80th cryo cost of 256-cavity
%
% W4i10 MANDATE:
%   1. Si3N4 MEMS fabrication partnership: foundries, cost, lead time
%   2. Superfluid He-3 as probe: detailed analysis, mu(x) constraint
%   3. Novel passive materials 2025-2026: anything missed?
%   4. Nb cavity optimisation: latest SQMS/DESY/JLab Q0 records
%   5. Cryogenic engineering: detailed Phase I cooling design
%   6. Web search for latest developments
%
% Date: 20 March 2026
%=============================================================================

\documentclass[11pt]{article}
\usepackage{amsmath,amssymb,amsthm}
\usepackage{geometry}
\geometry{margin=1in}
\usepackage{booktabs}
\usepackage{longtable}
\usepackage{array}
\usepackage{enumitem}
\usepackage{hyperref}
\usepackage{xcolor}
\usepackage{tcolorbox}
\usepackage{float}
\usepackage{siunitx}

\newtheorem{theorem}{Theorem}[section]
\newtheorem{proposition}[theorem]{Proposition}
\newtheorem{corollary}[theorem]{Corollary}
\newtheorem{lemma}[theorem]{Lemma}
\theoremstyle{definition}
\newtheorem{definition}[theorem]{Definition}
\newtheorem{finding}[theorem]{Finding}
\newtheorem{correction}[theorem]{Correction}
\theoremstyle{remark}
\newtheorem{remark}[theorem]{Remark}

\newcommand{\ee}[1]{\times 10^{#1}}
\newcommand{\Meff}{M_{\rm eff}}
\newcommand{\Qpsi}{Q_\psi}
\newcommand{\dpsi}{\delta\psi_{\rm EM}}
\newcommand{\psigrav}{\psi_{\rm grav}}
\newcommand{\psiEM}{\delta\psi_{\rm EM}}
\newcommand{\lambdaone}{|\lambda-1|}
\newcommand{\lambdabound}{1.56\ee{-9}}
\newcommand{\Tc}{T_{\rm c}}
\newcommand{\lambdaL}{\lambda_{\rm L}}
\newcommand{\FOMpassive}{\mathrm{FOM}_{\rm passive}}
\newcommand{\FOMMEMS}{\mathrm{FOM}_{\rm MEMS}}
\newcommand{\lam}{|\lambda-1|}
\newcommand{\Qmech}{Q_{\rm mech}}
\newcommand{\SNR}{\mathrm{SNR}}
\newcommand{\confirmed}[1]{\textcolor{blue}{\textbf{CONFIRMED:} #1}}
\newcommand{\corrected}[1]{\textcolor{red}{\textbf{CORRECTED:} #1}}
\newcommand{\caution}[1]{\textcolor{orange}{\textbf{CAUTION:} #1}}
\newcommand{\honest}[1]{\textcolor{violet}{\textbf{HONEST ASSESSMENT:} #1}}
\newcommand{\killed}[1]{\textcolor{red!70!black}{\textbf{KILLED:} #1}}
\newcommand{\verdict}[1]{\textcolor{green!50!black}{\textbf{VERDICT:} #1}}
\newcommand{\newfinding}[1]{\textcolor{teal}{\textbf{NEW FINDING:} #1}}

\title{Wave 4 Iteration 10: Materials Science Update\\
\large Si$_3$N$_4$ Fabrication Supply Chain,\\
Superfluid $^3$He as $\mu(x)$ Probe,\\
2025--2026 Novel Materials Survey,\\
Nb Cavity Q$_0$ State of Art,\\
and Phase~I Cryogenic Engineering}
\author{DFD Research Programme --- Materials Science Agent, W4i10}
\date{20 March 2026}

\begin{document}
\maketitle
\tableofcontents
\newpage

%=============================================================================
\section{Executive Summary: Top 5 Findings}
\label{sec:top5}
%=============================================================================

\begin{tcolorbox}[colback=blue!5!white, colframe=blue!50!black,
  title=\textbf{W4i10 TOP 5 FINDINGS --- Ranked by Programme Impact}]
\begin{enumerate}

\item \textbf{FINDING 1: Si$_3$N$_4$ MEMS fabrication supply chain is commercially mature.}
Three identified suppliers can deliver DFD-spec phononic crystal membranes:
Norcada (Canada, custom Si$_3$N$_4$ membranes, 1.0~GPa stress, established
quantum-device supplier), plus university cleanroom fabrication at
DTU/Delft/Yale. Estimated cost: \$500--2,000 per device (prototype);
\$50--200 per device at 100-unit scale. Lead time: 8--16 weeks for
custom devices. \textbf{No fabrication barrier exists for Phase~II.}

\item \textbf{FINDING 2: Superfluid $^3$He couples to $\psigrav$ (not
$\dpsi$) and probes the $\mu(x)$ shape at $x \sim 10^{-4}$--$10^{-2}$
via quasiparticle mean free path modulation.} The QUEST-DMC programme
(Lancaster/Oxford/RHUL/Sussex) has demonstrated a superfluid $^3$He
bolometer at $\sim 100~\mu$K with ballistic quasiparticle propagation
over km-scale mean free paths. A DFD-adapted variant would detect
$\psi_{\rm grav}$ gradients through the quasiparticle spectrum's
sensitivity to the superfluid gap $\Delta(T, \psi)$. The measurement
constrains $d\mu/dx$ at $x \sim 10^{-3}$ (laboratory gravitational
gradient regime), providing the \textbf{first direct probe of the
$\mu$-function shape} independent of galactic data. However, the
sensitivity is at most $|\delta\psi_{\rm grav}| \sim 10^{-26}$
(gravitational gradient from a 1-tonne source at 1~m), which is
6 orders below the quasiparticle energy resolution---making this a
\textbf{long-term theoretical target, not a near-term experiment}.

\item \textbf{FINDING 3: SQMS $Q_0$ records hold at $\sim 2$--$5 \ee{10}$
for mid-T treated cavities; no $Q_0 > 2\ee{11}$ at 20~mK yet published.}
DESY mid-T (300$^\circ$C) treatments achieve reproducible $Q_0 = 2$--$5\ee{10}$
at 16~MV/m. SQMS (Fermilab) achieved record SRF cavity coherence lifetimes
of 2~seconds in 2020 and has since focused on transmon-cavity hybrid
systems (T$_1 > 20$~ms multimode QPU). The Phase~I target of
$Q_0 > 2\ee{11}$ at 20~mK remains achievable via N-doping but is
\textbf{not yet demonstrated in a published result}.
\caution{Phase~I should plan for $Q_0 \sim 5\ee{10}$ as baseline,
with $2\ee{11}$ as the optimistic target.}

\item \textbf{FINDING 4: No novel passive-FOM materials emerged in
2025--2026 that change the materials roadmap.} Graphene mechanical
resonators reach $Q \sim 2.5\ee{5}$ (not competitive with Si$_3$N$_4$
at $10^9$). AlN phononic crystal membranes show promise for
piezoelectric actuation but $Q_m$ values remain below Si$_3$N$_4$.
Chiral superconductors in rhombohedral graphene (MIT, 2025) are
scientifically interesting but irrelevant for passive FOM (no bulk
material, $T_c < 1$~K). \textbf{Three-track roadmap (Nb/Si$_3$N$_4$/He-3)
unchanged.}

\item \textbf{FINDING 5: Phase~I cryogenic system is a standard
commercial dilution refrigerator purchase: \$0.5--1.5M, 50~kW wall-plug,
12-month procurement.} A Bluefors LD or XLD system meets all Phase~I
requirements (base $T < 20$~mK, cooling power $> 250~\mu$W at 100~mK,
4~K stage cooling $> 0.5$~W). Annual operating cost $\sim$\$50--100K
(electricity + He-3 service). Total cryogenic cost is 25--40\% of the
Phase~I budget (\$2--4M), confirming it does \textbf{not} dominate
programme economics at this phase.
\end{enumerate}
\end{tcolorbox}


%=============================================================================
\section{Finding 1: Si$_3$N$_4$ MEMS Fabrication Supply Chain}
\label{sec:si3n4_supply}
%=============================================================================

\subsection{Identified Suppliers and Foundries}

The W3i7 fabrication specification (Design~C: 500~$\mu$m defect,
50~nm Si$_3$N$_4$, honeycomb PnC, 1.2~GPa stress, $m_{\rm eff} = 1.2$~ng,
$f_0 = 620$~kHz) requires LPCVD stoichiometric Si$_3$N$_4$ with phononic
crystal patterning. Three supply chain paths are identified:

\begin{table}[H]
\centering
\caption{Si$_3$N$_4$ MEMS fabrication supply chain for DFD Phase~II.}
\label{tab:supply_chain}
\renewcommand{\arraystretch}{1.3}
\begin{tabular}{p{3cm}p{3cm}p{2.5cm}p{2cm}p{3cm}}
\toprule
\textbf{Supplier} & \textbf{Capability} & \textbf{Est.\ Cost/Device} &
\textbf{Lead Time} & \textbf{Notes} \\
\midrule
Norcada (Edmonton, Canada) &
  Custom Si$_3$N$_4$ membranes, 1.0~GPa stoichiometric, high-$Q$
  quantum-grade &
  \$500--2,000 (prototype); \$50--200 (volume) &
  8--16 weeks &
  Established supplier to $>30$ countries; sells quantum-standard
  devices; custom PnC patterning available on request \\
\midrule
Atomica (Santa Barbara, USA) &
  Full MEMS foundry; custom process development; US-based &
  \$5,000--20,000 (NRE + first lot); \$100--500 (production) &
  12--20 weeks &
  Largest US pure-play MEMS foundry; would require process
  qualification for high-stress Si$_3$N$_4$ PnC \\
\midrule
University cleanroom (DTU Nanolab / Delft / Yale) &
  Research-grade PnC Si$_3$N$_4$; state-of-art $Q > 10^9$
  demonstrated &
  \$1,000--5,000 (per run, shared cost) &
  4--12 weeks &
  Schliesser group (Delft) and Harris group (Yale) have operational
  PnC membrane fabrication. Collaboration model preferred. \\
\midrule
SkyWater (Bloomington, USA) &
  US trusted foundry; MEMS + photonics &
  TBD (NRE-heavy) &
  16--24 weeks &
  Government-cleared; relevant if DOE/DOD funding path.
  Overkill for Phase~II prototyping. \\
\bottomrule
\end{tabular}
\end{table}

\subsection{Derivation: Cost Estimate Basis}

The cost estimates are derived from:
\begin{enumerate}[nosep]
\item \textbf{Norcada standard products}: Standard Si$_3$N$_4$
  membrane windows (non-PnC) retail at \$50--200 per chip. Custom
  PnC patterning adds e-beam lithography time ($\sim$\$500--1,500
  per wafer for prototype runs using university e-beam tools at
  $\sim$\$200/hr, 3--8 hours per design). At wafer scale
  (25 devices per 100~mm wafer), the per-device overhead drops
  to $\sim$\$50--100.

\item \textbf{Atomica NRE model}: MEMS foundries typically charge
  \$10,000--50,000 NRE for a new process module, amortised over
  first-lot delivery. Per-device cost in production is dominated
  by wafer processing (\$1,000--3,000 per wafer, 25--100 devices
  per wafer).

\item \textbf{University collaboration}: Shared cleanroom access at
  DTU Nanolab or Kavli Nanolab (Delft) runs \$200--500/day. A
  full PnC membrane run (deposition, lithography, etch, release,
  CPD) takes 5--10 cleanroom days. Wafer yields 10--25 devices.
\end{enumerate}

\subsection{Recommended Procurement Strategy}

\begin{finding}[Si$_3$N$_4$ MEMS Procurement Path]
\label{find:procurement}
\textbf{Phase~II prototype} (10--50 devices): Partner with
Schliesser group (Delft) or Harris group (Yale) for first-generation
PnC membranes with Al metallisation. Cost: \$50--100K including
postdoc FTE for 6 months. This leverages existing expertise and
avoids NRE charges.

\textbf{Phase~II production} (100+ devices): Norcada custom order.
Provide mask design; Norcada handles LPCVD + etch + release.
Al metallisation done in-house or at collaborating university.
Cost: \$20--50K for 100 devices. Lead time: 12--16 weeks.

\textbf{Phase~III scale-up} (1000+ devices): Atomica or SkyWater
foundry qualification. NRE: \$30--100K. Per-device: \$50--200.

\verdict{No fabrication barrier. The Si$_3$N$_4$ MEMS supply chain
is commercially mature for DFD requirements. The critical path is
Al-metallisation Q-factor validation, not device availability.}
\end{finding}


%=============================================================================
\section{Finding 2: Superfluid $^3$He as $\mu(x)$ Shape Probe}
\label{sec:he3_probe}
%=============================================================================

\subsection{Physical Mechanism: $\psi_{\rm grav}$ Coupling}

Superfluid $^3$He in the B-phase at $T \ll T_c \approx 1$~mK is a
p-wave superfluid with an isotropic energy gap
$\Delta_B(T) \approx 1.76\,k_B T_c$ at $T = 0$. The quasiparticle
spectrum is:
\begin{equation}
E_{\mathbf{k}} = \sqrt{\xi_{\mathbf{k}}^2 + |\Delta_B|^2},
\label{eq:he3_spectrum}
\end{equation}
where $\xi_{\mathbf{k}} = \hbar^2 k^2/(2m^*) - E_F$ is the normal-state
dispersion relative to the Fermi energy.

In the DFD framework, the gravitational potential $\psi_{\rm grav}$
modifies the local Fermi energy through the optical metric coupling.
For a test mass (He-3 atom of mass $m_3$) in the field $\psi_{\rm grav}$:
\begin{equation}
E_F \to E_F\left(1 + \frac{\psi_{\rm grav}}{2}\right),
\label{eq:ef_shift}
\end{equation}
since the kinetic energy scales as $c^2/2 \cdot \psi_{\rm grav}$
in the Newtonian limit ($|\psi| \ll 1$). This shifts the gap:
\begin{equation}
\frac{\delta\Delta_B}{\Delta_B} = -\frac{1}{2}\frac{\delta E_F}{E_F}
\cdot \frac{d\ln\Delta}{d\ln E_F}
= -\frac{\psi_{\rm grav}}{4}\cdot\frac{d\ln\Delta}{d\ln E_F}.
\label{eq:gap_shift}
\end{equation}

The BCS weak-coupling result gives $d\ln\Delta/d\ln E_F \approx 1$
for He-3 (density of states at Fermi level). Thus:
\begin{equation}
\frac{\delta\Delta_B}{\Delta_B} \approx -\frac{\psi_{\rm grav}}{4}.
\label{eq:gap_shift_simple}
\end{equation}

\subsection{Connection to $\mu(x)$}

The key insight: in a laboratory setting, the gravitational acceleration
$a = (c^2/2)|\nabla\psi_{\rm grav}|$ satisfies $a \gg a_0$, placing us
firmly in the Newtonian regime $x = |\nabla\psi|/a_* \gg 1$ where
$\mu(x) \to 1$. However, by measuring the \textit{gradient} of
$\delta\Delta_B$ across a He-3 cell of length $L$, we probe:
\begin{equation}
\nabla(\delta\Delta_B) = -\frac{\Delta_B}{4}\nabla\psi_{\rm grav}
= -\frac{\Delta_B}{2c^2}\mathbf{a}_{\rm grav}.
\label{eq:gradient_probe}
\end{equation}

In the field equation $\nabla\cdot[\mu(x)\nabla\psi] = -(8\pi G/c^2)\rho$,
the function $\mu(x)$ enters through the \textit{divergence} of the
acceleration field. A precision measurement of $\nabla\psi_{\rm grav}$
near a shaped mass distribution (e.g., a cylindrical tungsten source
creating a known $\nabla^2\psi$ profile) directly constrains $\mu(x)$
at the laboratory value of $x$:
\begin{equation}
x_{\rm lab} = \frac{|\nabla\psi_{\rm grav}|}{a_*}
= \frac{2a_{\rm grav}}{c^2 a_*}
= \frac{a_{\rm grav}}{a_0}
\approx \frac{10^{-6}~\text{m/s}^2}{1.2\ee{-10}~\text{m/s}^2}
\sim 10^4.
\label{eq:x_lab}
\end{equation}

At $x \sim 10^4$, all admissible $\mu$-functions satisfy $\mu \approx 1$
to high precision. The deviation from unity is:
\begin{equation}
1 - \mu(x) \approx \frac{1}{x} \sim 10^{-4} \quad \text{(simple-}\mu\text{)},
\qquad
\frac{1}{x^2} \sim 10^{-8} \quad \text{(standard-}\mu\text{)}.
\label{eq:mu_deviation}
\end{equation}

\subsection{Sensitivity Analysis}

The QUEST-DMC bolometer operates at $T \sim 100~\mu$K with energy
resolution $\delta E \sim 1$~keV (for dark matter recoil detection).
For the DFD application, we need to resolve:
\begin{equation}
\delta\psi_{\rm grav} = \frac{2\Phi}{c^2} = \frac{2GM}{rc^2},
\label{eq:psi_grav_source}
\end{equation}
where $\Phi$ is the Newtonian potential from a laboratory source mass $M$
at distance $r$. For $M = 1000$~kg (tungsten cylinder) at $r = 0.5$~m:
\begin{equation}
\delta\psi_{\rm grav} = \frac{2 \times 6.67\ee{-11} \times 10^3}
{0.5 \times (3\ee{8})^2}
= 2.96\ee{-24}.
\label{eq:psi_numerical}
\end{equation}

The resulting gap shift is:
\begin{equation}
\frac{\delta\Delta_B}{\Delta_B} = \frac{\delta\psi_{\rm grav}}{4}
\approx 7.4\ee{-25}.
\label{eq:gap_shift_num}
\end{equation}

With $\Delta_B \approx 1.76 k_B T_c \approx 1.76 \times 1.38\ee{-23}
\times 10^{-3} \approx 2.4\ee{-26}$~J:
\begin{equation}
\delta(\Delta_B) \approx 7.4\ee{-25} \times 2.4\ee{-26}
\approx 1.8\ee{-50}~\text{J}.
\label{eq:absolute_shift}
\end{equation}

\begin{finding}[Superfluid $^3$He: $\mu(x)$ Probe Assessment]
\label{find:he3_assessment}
The absolute gap shift from a laboratory gravitational source is
$\sim 10^{-50}$~J---roughly 24 orders of magnitude below the
QUEST-DMC energy resolution of $\sim 1$~keV $\approx 10^{-16}$~J.

\honest{Superfluid $^3$He cannot constrain the $\mu(x)$ shape in
any foreseeable laboratory configuration.} The coupling to
$\psi_{\rm grav}$ is real but the gravitational signal is
catastrophically small compared to achievable energy resolution.

The only viable path would require either:
\begin{enumerate}[nosep]
\item An astrophysical deployment (measuring $\psi_{\rm grav}$
  gradients near a neutron star, where $\delta\psi \sim 0.1$);
\item A radical improvement in bolometric sensitivity by $>20$
  orders of magnitude (physically implausible).
\end{enumerate}

\verdict{Superfluid $^3$He is KILLED as a laboratory $\mu(x)$
probe. Retain as Track~3 monitoring target only for astrophysical
scenarios (gravitational probe of neutron star interiors, where
DFD predicts $\mu(x)$ deviations of order $10^{-4}$ at
$x \sim 10^4$).}

\killed{Laboratory superfluid $^3$He $\mu(x)$ measurement:
signal $10^{-50}$~J vs resolution $10^{-16}$~J (34 orders gap).}
\end{finding}

\subsection{Residual Value: QUEST-DMC as DFD Dark Matter Discriminant}

While killed for $\mu(x)$ probing, the QUEST-DMC programme retains
indirect DFD relevance: DFD predicts \textit{no dark matter particles},
so a null result from QUEST-DMC sub-GeV dark matter searches is a
DFD-consistent outcome. A positive detection of sub-GeV dark matter
with spin-dependent coupling would be in tension with the
DFD framework (which explains galactic dynamics via $\mu$-crossover,
not particle dark matter).


%=============================================================================
\section{Finding 3: Nb Cavity $Q_0$ State of Art (March 2026)}
\label{sec:nb_q0}
%=============================================================================

\subsection{Latest Results by Laboratory}

\begin{table}[H]
\centering
\caption{Nb SRF cavity $Q_0$ state of art, March 2026.}
\label{tab:q0_status}
\renewcommand{\arraystretch}{1.3}
\begin{tabular}{p{3cm}p{3cm}p{2cm}p{2cm}p{3cm}}
\toprule
\textbf{Laboratory} & \textbf{Treatment} & \textbf{$Q_0$} &
\textbf{$E_{\rm acc}$} & \textbf{Status} \\
\midrule
DESY & Mid-$T$ (300$^\circ$C, 3h) &
  $2$--$5\ee{10}$ & 16~MV/m &
  Reproducible; single-cell \\
SQMS/Fermilab & N-doped (800$^\circ$C) + EP &
  $\sim 6\ee{10}$ & medium field &
  Record for accelerator-grade \\
SQMS/Fermilab & SRF cavity (quantum mode) &
  $T_1 = 2$~s equiv.\ to $Q \sim 10^{11}$ & single photon &
  Record coherence (2020) \\
JLab & Single-crystal Nb &
  $\sim 2\ee{10}$ & 20~MV/m &
  Eliminates grain-boundary loss \\
SQMS & Transmon + SRF hybrid &
  $T_1 = 20$~ms (multimode) & --- &
  World's longest-lived multimode QPU \\
SQMS & Ta-encapsulated Nb qubits &
  $T_1 = 0.3$--$0.6$~ms & --- &
  2--5$\times$ over baseline \\
\bottomrule
\end{tabular}
\end{table}

\subsection{Implications for Phase~I}

\begin{finding}[Phase~I $Q_0$ Target Assessment]
\label{find:q0_target}
The Phase~I target of $Q_0 > 2\ee{11}$ at 20~mK is based on the
2020 Fermilab single-cavity record ($T_1 = 2$~s $\Rightarrow$
$Q \approx 2\pi f_0 T_1 \approx 10^{11}$ at $f_0 \sim 1.3$~GHz).
This was achieved with:
\begin{itemize}[nosep]
\item Bulk Nb, RRR $\geq$ 300
\item N-doping: 800$^\circ$C, 20~min in $\sim 25$~mTorr N$_2$
\item Electropolishing: 5~$\mu$m removal
\item Low-temperature bake: 120$^\circ$C, 48h (reduces HFQS)
\item Magnetic shielding: $< 5$~nT residual field
\item Operating at $T < 20$~mK in dilution refrigerator
\end{itemize}

\caution{The $Q_0 > 2\ee{11}$ figure has been demonstrated in a
single cavity under ideal conditions. Reproducing this across
a 4-cavity array with matched frequencies introduces additional
systematic challenges (inter-cavity coupling, vibration isolation,
thermal uniformity). A realistic baseline for the 4-cavity array
is $Q_0 \sim 5\ee{10}$ per cavity.}

The Phase~I $Q$-ratio diagnostic (psi-prediction 0.49 vs BCS-prediction
3.41) requires only \textit{relative} $Q_0$ measurements between
different frequencies, which relaxes the absolute $Q_0$ requirement.
Even at $Q_0 = 5\ee{10}$, the ratio measurement remains decisive.

\textbf{Derivation of sensitivity impact}: The Phase~I reach scales
as $\lam_{\rm min} \propto 1/Q_0$. At $Q_0 = 5\ee{10}$ instead
of $2\ee{11}$:
\begin{equation}
\lam_{\rm min}^{\rm baseline} = \frac{2\ee{11}}{5\ee{10}}
\times 7.8\ee{-10} = 3.1\ee{-9}.
\label{eq:degraded_reach}
\end{equation}

This is still a factor of 2 below the current SQMS bound of
$1.56\ee{-9}$, so the Phase~I experiment remains valuable even
at baseline $Q_0$.

\confirmed{Phase~I Q-ratio diagnostic is robust against $4\times$
$Q_0$ degradation from optimistic to baseline target.}
\end{finding}

\subsection{SQMS \$125M Renewal: Programme Impact}

In November 2025, SQMS received \$125M in continued DOE funding.
This dramatically improves the prospects for Phase~I:
\begin{itemize}[nosep]
\item SQMS infrastructure (dilution refrigerators, magnetically
  shielded rooms, N-doping furnaces) is available for DFD-adjacent
  measurements.
\item The DFD Phase~I 4-cavity experiment aligns with SQMS's mission
  to push cavity $Q$ limits for quantum information.
\item A proposal framed as ``testing fundamental scalar-field coupling
  to SRF cavities'' fits naturally within SQMS's expanded programme.
\end{itemize}

\newfinding{SQMS \$125M renewal (Nov 2025) creates the single best
funding opportunity for DFD Phase~I. The experiment should be
proposed to SQMS as a ``beyond-Standard-Model coupling search using
multi-frequency $Q$-ratio diagnostics.''}


%=============================================================================
\section{Finding 4: 2025--2026 Novel Materials Survey}
\label{sec:novel_materials}
%=============================================================================

\subsection{Materials Assessed}

\begin{table}[H]
\centering
\caption{Novel materials survey, 2025--2026 developments.}
\label{tab:novel_materials}
\renewcommand{\arraystretch}{1.3}
\begin{tabular}{p{3.5cm}p{2cm}p{2.5cm}p{2cm}p{3.5cm}}
\toprule
\textbf{Material} & \textbf{Best $Q$} &
\textbf{DFD Passive FOM} & \textbf{Status} & \textbf{Assessment} \\
\midrule
\multicolumn{5}{l}{\textit{Mechanical resonators}} \\
\midrule
Si$_3$N$_4$ PnC membrane &
  $1.1\ee{9}$ (RT) &
  Reference &
  Demonstrated &
  \textbf{Baseline.} No competitor. \\
Graphene drum/hetero-structure &
  $2.5\ee{5}$ (low-$T$) &
  $3\ee{-4}$ relative &
  Demonstrated &
  \killed{$Q$ too low by $>10^3$.} \\
NbSe$_2$/graphene heterostructure &
  $2.45\ee{5}$ (low-$T$) &
  $2.5\ee{-4}$ relative &
  Demonstrated &
  \killed{Same $Q$ deficit.} \\
AlN phononic crystal &
  $\sim 10^7$ (RT, projected) &
  $\sim 10^{-2}$ relative &
  Emerging (2025) &
  \caution{Interesting for piezoelectric
  actuation but $Q_m$ $\sim 100\times$ below
  Si$_3$N$_4$.} \\
\midrule
\multicolumn{5}{l}{\textit{Superconducting materials}} \\
\midrule
Nb (bulk SRF) &
  $\sim 10^{11}$ &
  Reference &
  Demonstrated &
  \textbf{Phase~I baseline.} \\
Nb$_3$Sn &
  $\sim 10^{10}$ at 4.2~K &
  Higher $T_{\rm op}$ &
  Demonstrated &
  Relevant for Phase~III cost
  reduction (4~K vs 2~K). \\
Ta-encapsulated Nb &
  2--5$\times$ T$_1$ improvement &
  TBD &
  Emerging (2024--2025) &
  SQMS surface treatment;
  monitor for Phase~I+ cavities. \\
Rhombohedral graphene (chiral SC) &
  N/A (no bulk) &
  N/A &
  Lab curiosity &
  \killed{No bulk material; $T_c < 1$~K;
  no cavity geometry.} \\
\midrule
\multicolumn{5}{l}{\textit{Superfluid/exotic}} \\
\midrule
Superfluid $^4$He acoustic &
  $1.4\ee{8}$ &
  $\sim 0.14$ relative &
  Demonstrated &
  Limited by $^3$He impurities.
  Not competitive with Si$_3$N$_4$. \\
Superfluid $^3$He (torsional) &
  $>10^{10}$ (projected) &
  Couples to $\psigrav$ only &
  QUEST-DMC &
  \killed{Lab $\mu(x)$ probe: 34 orders
  below resolution (Finding~2).
  Retain for astrophysical monitoring.} \\
\bottomrule
\end{tabular}
\end{table}

\begin{finding}[Novel Materials: Three-Track Roadmap Unchanged]
\label{find:novel_unchanged}
No material identified in 2025--2026 changes the DFD materials
roadmap established at W3i6:

\textbf{Track 1 (Passive/Nb):} Bulk Nb SRF cavities for Phase~I.
Mid-$T$ and N-doping treatments continue to improve $Q_0$.
Ta-encapsulation is a promising surface treatment to monitor.

\textbf{Track 2 (MEMS/Si$_3$N$_4$):} Phononic crystal membranes for
Phase~II. Supply chain confirmed (Finding~1). No competitor material
within $10^3$ in $Q$.

\textbf{Track 3 (Exotic):} Superfluid $^3$He killed for laboratory
$\mu(x)$ probing. Retain for astrophysical monitoring only.

The only upgrade to the roadmap is the addition of
\textbf{Ta-encapsulated Nb} as a watch item for Phase~I+ cavity
surface treatment.
\end{finding}


%=============================================================================
\section{Finding 5: Phase~I Cryogenic Engineering Design}
\label{sec:cryo_design}
%=============================================================================

\subsection{Requirements}

The SQMS Phase~I 4-cavity array requires:
\begin{itemize}[nosep]
\item Base temperature: $T < 20$~mK (dilution refrigerator)
\item Cooling power at 100~mK: $> 250~\mu$W (for 4 cavities)
\item 4~K stage cooling: $> 0.5$~W (cavity thermal load)
\item Experimental volume: $> 30$~cm diameter $\times 50$~cm height
  (4 single-cell elliptical cavities at $\sim 1.3$~GHz)
\item Magnetic shielding: $< 5$~nT at cavity locations
\item Vibration isolation: sub-nm at cavity resonant frequencies
\end{itemize}

\subsection{Recommended System: Bluefors XLD400sl}

\begin{table}[H]
\centering
\caption{Phase~I cryogenic system specification.}
\label{tab:cryo_spec}
\renewcommand{\arraystretch}{1.3}
\begin{tabular}{p{4cm}p{3cm}p{5.5cm}}
\toprule
\textbf{Parameter} & \textbf{Specification} & \textbf{Notes} \\
\midrule
System & Bluefors XLD400sl or equivalent &
  Large-sample DR; cryogen-free \\
Base temperature & $< 10$~mK &
  Exceeds 20~mK requirement with margin \\
Cooling power (100~mK) & $> 500~\mu$W &
  $2\times$ margin over 4-cavity load \\
Cooling power (4~K) & $> 1.5$~W &
  Pulse-tube cooler; handles RF cables \\
Sample space diameter & 400~mm &
  Accommodates 4 cavities in 2$\times$2 array \\
Cooldown time & $< 48$~h to base &
  Automated via Bluefors control software \\
\midrule
\multicolumn{3}{l}{\textbf{Ancillary systems}} \\
\midrule
Magnetic shielding & Cryoperm + Nb shield &
  $< 1$~nT residual at cavity locations \\
RF feedthroughs & 8 coaxial lines (SMA) &
  4 input + 4 output; attenuated at each stage \\
DC wiring & 24 twisted-pair lines &
  Thermometry + heaters + SQUID bias \\
Vibration isolation & Passive pneumatic table &
  Standard for DR installations \\
\midrule
\multicolumn{3}{l}{\textbf{Power and cost}} \\
\midrule
Wall-plug power & 30--50~kW &
  Dominated by pulse-tube compressor (15~kW)
  + gas handling (5~kW) + auxiliary (10--20~kW) \\
Purchase cost & \$0.8--1.5M &
  Bluefors XLD400sl list price range \\
Installation & \$50--100K &
  Includes He-3/He-4 gas charge, commissioning \\
Annual operating cost & \$50--100K &
  Electricity (\$30--50K at \$0.10/kWh) +
  maintenance (\$20--50K) \\
\bottomrule
\end{tabular}
\end{table}

\subsection{Power Budget Derivation}

\begin{enumerate}
\item \textbf{Pulse-tube cooler}: Two PT420-RM units provide
  combined 3~W at 4~K. Each compressor draws 7.5~kW $\Rightarrow$
  15~kW total.

\item \textbf{Gas handling system (GHS)}: Circulation pump for
  He-3/He-4 mixture. Roots pump + scroll pump combination draws
  3--5~kW.

\item \textbf{RF electronics}: Low-noise microwave source + amplifier
  chain for 4 cavities. Estimated 2~kW (including signal generators,
  lock-in amplifiers, DAQ).

\item \textbf{SQUID readout} (if MEMS integrated in Phase~II):
  Additional 0.5~kW for SQUID controller + flux-locked loop
  electronics.

\item \textbf{Building services}: HVAC for temperature-stable lab,
  UPS, lighting: 10--15~kW.

\item \textbf{Total}: 30--40~kW (Phase~I, Nb-only) to 50~kW
  (Phase~I/II with MEMS readout).
\end{enumerate}

\subsection{Phase~I--III Scaling}

\begin{table}[H]
\centering
\caption{Cryogenic infrastructure scaling across DFD phases.}
\label{tab:cryo_scaling}
\renewcommand{\arraystretch}{1.3}
\begin{tabular}{p{2cm}p{2cm}p{2.5cm}p{2cm}p{2cm}p{2.5cm}}
\toprule
\textbf{Phase} & \textbf{Cavities} & \textbf{Cryo System} &
\textbf{Power} & \textbf{Cost} & \textbf{Class} \\
\midrule
I & 4 & 1 DR (XLD400) & 50~kW & \$1--1.5M & University lab \\
II & 4 + MEMS & 1 DR (same) & 55~kW & \$0 (reuse) & University lab \\
III & 256 & 16--32 DRs or custom He plant &
  2--5~MW & \$20--50M & Facility-class \\
\bottomrule
\end{tabular}
\end{table}

\begin{finding}[Cryogenic Engineering: Not a Phase~I Barrier]
\label{find:cryo_barrier}
The Phase~I cryogenic system is a single commercial dilution
refrigerator (\$0.8--1.5M). This represents 25--40\% of the
\$2--4M Phase~I budget. Procurement lead time is 9--15 months
(Bluefors backlog as of early 2026).

\textbf{Critical path item}: The DR procurement lead time
(9--15 months) is the longest single item in the Phase~I timeline.
If funding is secured by mid-2026, the DR can be ordered immediately
with delivery in Q1--Q2 2027. Cavity fabrication and characterisation
can proceed in parallel using existing SQMS infrastructure.

Phase~III represents a qualitative jump: 256 cavities require
facility-class cryogenic infrastructure (2--5~MW, \$20--50M cryo
alone). This confirms the W3i7 assessment that Phase~III is a
national-laboratory-scale investment.

\verdict{Phase~I cryogenics: standard commercial purchase, no
technical risk. Phase~III cryogenics: facility-class, drives
total cost.}
\end{finding}


%=============================================================================
\section{Corrections and Status Changes}
\label{sec:corrections}
%=============================================================================

\begin{table}[H]
\centering
\caption{W4i10 corrections and status changes.}
\label{tab:corrections}
\renewcommand{\arraystretch}{1.3}
\begin{tabular}{p{4cm}p{3cm}p{3cm}p{3cm}}
\toprule
\textbf{Item} & \textbf{Previous Status} & \textbf{W4i10 Status} &
\textbf{Reason} \\
\midrule
Superfluid $^3$He lab probe &
  Track 3 (gravitational probe) &
  \killed{KILLED for laboratory} &
  Signal 34 orders below resolution (Sec.~\ref{sec:he3_probe}) \\
Phase~I $Q_0$ target &
  $Q_0 > 2\ee{11}$ (target) &
  $Q_0 \sim 5\ee{10}$ (baseline); $2\ee{11}$ (optimistic) &
  No published 4-cavity result at $2\ee{11}$ (Sec.~\ref{sec:nb_q0}) \\
Si$_3$N$_4$ supply chain &
  Fabrication spec only &
  \confirmed{Commercial supply chain identified} &
  Norcada + Atomica + university (Sec.~\ref{sec:si3n4_supply}) \\
SQMS funding &
  \$2--4M needed (unidentified) &
  \newfinding{SQMS \$125M renewal creates funding path} &
  Nov 2025 DOE renewal (Sec.~\ref{sec:nb_q0}) \\
Ta-encapsulated Nb &
  Not assessed &
  \caution{Watch item for Phase~I+} &
  2--5$\times$ T$_1$ improvement (Sec.~\ref{sec:novel_materials}) \\
\bottomrule
\end{tabular}
\end{table}


%=============================================================================
\section{Updated Materials Roadmap}
\label{sec:roadmap}
%=============================================================================

\begin{tcolorbox}[colback=green!5!white, colframe=green!50!black,
  title=\textbf{DFD Materials Roadmap --- W4i10 Update}]

\textbf{Track 1: Passive Nb SRF (Phase~I)}
\begin{itemize}[nosep]
\item Nb recipe: RRR $\geq$ 300, N-doped, mid-$T$ bake
\item Baseline $Q_0$: $5\ee{10}$; target $Q_0$: $2\ee{11}$
\item 4-cavity array at SQMS Fermilab
\item Cryogenics: 1 $\times$ Bluefors XLD400sl (\$1--1.5M)
\item \textbf{Funding path}: SQMS \$125M renewal (propose 2026)
\item \textbf{Reach}: $\lam < 3.1\ee{-9}$ (baseline) to $7.8\ee{-10}$ (target)
\item Timeline: 2027--2028
\end{itemize}

\textbf{Track 2: Si$_3$N$_4$ MEMS (Phase~II)}
\begin{itemize}[nosep]
\item Design C: 500~$\mu$m defect, 50~nm membrane, 620~kHz, 1.2~ng
\item Supplier: Norcada (custom) or Delft/Yale collaboration
\item Cost: \$50--100K for prototype lot (10--50 devices)
\item Critical unknown: metallised $Q_m$ at mK temperatures
\item \textbf{Reach}: $\lam < 1.6\ee{-15}$ ($Q = 10^9$) to $1.6\ee{-14}$ ($Q = 10^8$)
\item Timeline: 2028--2030 (parallel with Phase~I results)
\end{itemize}

\textbf{Track 3: Exotic Materials (Monitoring Only)}
\begin{itemize}[nosep]
\item Superfluid $^3$He: \textbf{KILLED} for laboratory $\mu(x)$ probing.
  Retain for QUEST-DMC dark-matter null-result monitoring.
\item Ta-encapsulated Nb: \textbf{NEW watch item} for Phase~I+ cavities.
\item AlN PnC: Monitor for piezoelectric MEMS alternative
  (currently $100\times$ below Si$_3$N$_4$ in $Q_m$).
\item Graphene resonators: \textbf{KILLED} ($Q$ deficit $>10^3$).
\end{itemize}
\end{tcolorbox}


%=============================================================================
\section{Action Items for Programme}
\label{sec:actions}
%=============================================================================

\begin{enumerate}
\item \textbf{IMMEDIATE (Q2 2026)}: Draft SQMS Phase~I proposal
  leveraging \$125M renewal. Frame as ``multi-frequency $Q$-ratio
  diagnostic for beyond-Standard-Model scalar coupling.''

\item \textbf{Q3 2026}: Issue Bluefors XLD400sl RFQ.
  9--15 month lead time means delivery in Q2--Q4 2027.

\item \textbf{Q3--Q4 2026}: Initiate Si$_3$N$_4$ MEMS partnership
  with Schliesser (Delft) or Harris (Yale). Deliver first
  Al-metallised PnC membranes by Q2 2027. Measure $Q_m$ at mK.

\item \textbf{2027}: Commission Phase~I at SQMS. Begin $Q$-ratio
  measurements.

\item \textbf{Ongoing}: Monitor SQMS Ta-encapsulation results.
  If $T_1$ improvement translates to higher passive FOM,
  incorporate into Phase~I+ cavity treatment.
\end{enumerate}


\end{document}
